\documentclass[11pt,a4paper]{article}
\usepackage[T1]{fontenc}
\usepackage{lmodern,amsmath,amssymb,amsthm,microtype}
\usepackage[margin=24mm]{geometry}
\usepackage{xcolor}
\usepackage[hidelinks]{hyperref}
\definecolor{ink}{HTML}{17354A}
\hypersetup{pdftitle={A finiteness result for y cubed plus y equals x to the fourth plus x},pdfauthor={}}
\setlength{\parindent}{0pt}
\setlength{\parskip}{5pt}
\setlength{\abovedisplayskip}{8pt plus 2pt minus 2pt}
\setlength{\belowdisplayskip}{8pt plus 2pt minus 2pt}
\newtheorem{theorem}{Theorem}
\newtheorem*{proposition}{Elementary proposition}
\newcommand{\Q}{\mathbb{Q}}
\newcommand{\C}{\mathbb{C}}
\newcommand{\Z}{\mathbb{Z}}
\newcommand{\PP}{\mathbb{P}}
\begin{document}
\begin{center}
{\Large\bfseries\color{ink} A finiteness result for \(y^3+y=x^4+x\)}\\[5pt]
{\small A correction to the proposed infinitude assertion}
\end{center}

\textbf{Status of the answer.}
The assertion of infinitely many integer solutions is false. The proof below
uses Siegel's theorem and the genus formula for a smooth plane curve; their
statements are supplied, but they are not proved. Thus the argument is
\emph{unconditional}, but it does \emph{not} meet the requested fully
self-contained standard. No complete enumeration of the integer solutions is
claimed. No internet search was used.

\begin{theorem}
There are only finitely many pairs \((x,y)\in\Z^2\) satisfying
\begin{equation}\label{eq:main}
 y^3+y=x^4+x.
\end{equation}
\end{theorem}

\begin{proof}
We verify explicitly the geometric hypotheses of the finiteness theorem.
Homogenizing \eqref{eq:main} gives the projective plane curve
\[
 C:\quad F(X,Y,Z)=Y^3Z+YZ^3-X^4-XZ^3=0.
\]
All smoothness and irreducibility checks below are over \(\C\).

\textbf{1. The point at infinity is unique and smooth.}
Setting \(Z=0\) gives \(X=0\), so the only point at infinity is
\(P=[0:1:0]\). Since
\[
 F_Z=Y^3+3YZ^2-3XZ^2,\qquad F_Z(P)=1,
\]
the gradient is nonzero at \(P\).

\textbf{2. There are no affine singularities.}
For \(f(x,y)=y^3+y-x^4-x\), a singular point would satisfy
\[
 f=0,\qquad f_x=-4x^3-1=0,\qquad f_y=3y^2+1=0.
\]
Consequently \(x^3=-1/4\) and \(y^2=-1/3\). Substitution into \(f=0\)
then gives
\[
 \frac23y=\frac34x,\qquad
 y=\frac98x,\qquad x^2=-\frac{64}{243}.
\]
But these equations also force
\[
 x=\frac{x^3}{x^2}=\frac{243}{256},
\]
whose square is positive, contradicting \(x^2=-64/243\).
Thus \(C\) is smooth everywhere.

\textbf{3. The curve is geometrically irreducible.}
Suppose \(F=GH\), with nonconstant homogeneous complex polynomials of
degrees \(d,e>0\). Since \(F(X,Y,0)=-X^4\), their restrictions to \(Z=0\)
must be \(cX^d\) and \(c'X^e\), where \(cc'=-1\). Hence both factors
vanish at \(P\), giving
\(\nabla F(P)=H(P)\nabla G(P)+G(P)\nabla H(P)=0\).
This contradicts step 1 and proves irreducibility.

\textbf{4. Apply the genus formula and Siegel's theorem.}
The genus formula states that a smooth geometrically irreducible projective
plane curve of degree \(d\) has genus \((d-1)(d-2)/2\). Here
\[
 g(C)=\frac{(4-1)(4-2)}2=3.
\]
The form of \textbf{Siegel's theorem} needed here is: an affine
geometrically irreducible curve over \(\Q\) whose smooth projective
completion has positive genus has only finitely many integer points in any
fixed affine embedding. Applying it to \(C\setminus\{P\}\), with coordinates
\((x,y)\), proves the theorem.
\end{proof}

\newpage
\begin{center}
{\large\bfseries\color{ink} Intuition and elementary information}
\end{center}

The leading terms suggest trying
\[
 x=t^3,\qquad y=t^4,
\]
because then \(y^3=x^4=t^{12}\). In the actual equation this leaves
\(t^4=t^3\), so the trial works only for \(t=0,1\). The lower-order terms
therefore matter even though they are small compared with the leading terms.
This failed trial alone would not rule out some other infinite family.

The geometric calculation resolves that difficulty. The simpler curve
\(y^3=x^4\) has a singularity at the origin and admits the displayed
parametrization. For the curve in \eqref{eq:main}, the derivative calculation
shows that no singularity remains, including over the complex numbers. Its
smooth quartic structure gives genus \(3\), and Siegel's theorem supplies the
arithmetic finiteness conclusion. Merely comparing the degrees \(3\) and \(4\)
would not establish finiteness.

Some useful information can also be proved entirely elementarily. Put
\(h(u)=u^3+u\). For real \(v>u\),
\[
 h(v)-h(u)=(v-u)(v^2+uv+u^2+1)>0,
\]
so \(h\) is strictly increasing. For every integer \(x\),
\(x^4+x\geq0\); hence any integer solution has \(y\geq0\). At
\(x=-1,0,1\), strict monotonicity determines exactly the solutions
\[
 \boxed{(-1,0),\qquad(0,0),\qquad(1,1).}
\]
For \(a=|x|\geq2\),
\[
 x^4+x-h(a)\geq a^4-a-a^3-a
 =a\bigl(a^2(a-1)-2\bigr)>0,
\]
so every further solution would satisfy \(y>|x|\).

\begin{proposition}
The only integer solutions of \eqref{eq:main} for which \(x\) is an integer
cube are \((-1,0),(0,0),(1,1)\).
\end{proposition}
\begin{proof}
Write \(x=t^3\). The cases \(t=-1,0,1\) were handled above. Suppose
\(|t|\geq2\), and put \(q=t^4\geq16\). Then
\[
 x^4+x=q^3+t^3<q^3+q=h(q),
\]
because \(t^3<q\). On the other hand, \(|t|^3\leq q\) gives
\begin{align*}
 q^3+t^3-h(q-1)
 &=3q^2-4q+2+t^3\\
 &\geq3q^2-5q+2=(3q-2)(q-1)>0.
\end{align*}
Thus \(h(q-1)<x^4+x<h(q)\). Strict monotonicity makes
\(h(y)=x^4+x\) impossible for integral \(y\), since no integer lies
strictly between \(q-1\) and \(q\).
\end{proof}

The proposition settles every cube value of \(x\), but does not settle the
noncube values. The global result proved on the first page is finiteness;
it supplies neither an explicit bound for \(|x|\) nor a proof that the three
displayed solutions are the complete list. The reliance on Siegel's theorem
is the substantive gap between this argument and a fully self-contained
treatment.

\end{document}

